% %%%%%%%%%%%%%%%%%%%%%%%%%%%%%%%%%%%%%%%%%%%%%%%%%%%%%%%%%%%%%%%%%%%%%%%%%%%%%%%%%%%%%%%
% Introduction
% %%%%%%%%%%%%%%%%%%%%%%%%%%%%%%%%%%%%%%%%%%%%%%%%%%%%%%%%%%%%%%%%%%%%%%%%%%%%%%%%%%%%%%%

\FancyHeadRight{\sectitle}

\renewcommand{\sectitle}{Introduction}

\phantomsection
\chapter*{\sectitle}
\addstarredchapter{\sectitle}

\lettrine{T}{he} rapid growth of data-intensive applications has made sparse and \gls{irrData} structures increasingly prevalent in scientific computing, machine learning, and graph analytics.
However, efficiently processing sparse matrices remains a fundamental challenge due to their inherent irregularity, which leads to unpredictable memory access patterns, limited opportunities for data reuse, and suboptimal utilization of hardware resources.
Existing hardware accelerators have demonstrated various solutions, yet they often rely on restrictive assumptions, exhibit limited scalability, or fail to fully exploit the memory bandwidth.

This thesis addresses the following central question:

\bigskip
\begin{highlight}{Problem Statement}{Mulberry}
    \colorbf{Mulberry}{\emph{How can hardware and architectural mechanisms be designed to efficiently support computation with sparse data, while ensuring scalability, adaptability, and high performance across diverse workloads?}}
\end{highlight}
\bigskip

To tackle this problem, we investigate the limitations of existing algorithms and accelerators, propose a novel cache architecture with support for read-accumulate-write operations, tailored to sparse data.
We then evaluate the behaviour of this specialized cache through modeling, simulation, and hardware prototyping.

\bigskip\bigskip

The manuscript is structured as follows.
We begin with an in-depth discussion of \gls{irrData}, with a particular emphasis on sparse matrices (Chapter~\ref{cha:background}).
This includes a study of fundamental sparse matrix multiplication kernels and their associated algorithms, highlighting the main hardware challenges that arise when dealing with sparse data.
A \gls{highLevel} analysis of algorithms for sparse data computing is provided, with a focus on basic storage schemes and opportunities for data reuse.
Building on this foundation, we review the state-of-the-art in hardware accelerators designed for efficient sparse matrix computing (Chapter~\ref{cha:soa}).
This review identifies common characteristics, proposes a classification of existing accelerators and design tendencies, offers a comparative analysis, and identifies the major tendencies across designs.
This analysis guided our subsequent architectural choices.

Following this survey, we introduce a novel cache architecture specifically designed to address the challenge of read-accumulate-write operations on sparse data, called ``random'' \glspl{reduction} (Chapter~\ref{cha:spdc}).
The proposed cache integrates \gls{reduction} support and organizes its storage into multiple regions optimized for denser and sparser data regions.
A Python model was developed to simulate memory transactions in this cache and to compare its behavior with that of a generic cache.
To further investigate its performance, we complement this analysis with a probabilistic model, that enables rapid simulation and design parameters exploration depending on matrix structure (Chapter~\ref{cha:theory}).

At the micro-architectural level, we detail the virtual partitioning of cache memory into multiple configurable regions, which ensures both flexibility and adaptability to diverse workloads (Chapter~\ref{cha:arch}).
The proposed cache is then evaluated through \acrshort{rtl} simulation, which reveals two distinct classes of matrices that exhibit different performance profiles.
Finally, the manuscript outlines potential future enhancements to improve both performance and scalability, including a multicore extension of the cache architecture and optimizations targeting better usability (Chapter~\ref{cha:optim}).

\section*{Contributions of the Thesis}

The main contributions of this thesis can be summarized as follows:

\begin{enumerate}
    \item \textbf{Analysis of sparse data computing} {\raggedleft \hfill (Chapter~\ref{cha:background})}
    \begin{itemize}
        \item Study of fundamental sparse matrix multiplication kernels and algorithms.
        \item Identification of the hardware challenges associated with sparse data processing.
        \item \Gls{highLevel} analysis of algorithms for sparse data computing, with emphasis on storage schemes and data reuse opportunities.
    \end{itemize}

    \item \textbf{Survey and classification of state-of-the-art accelerators} {\raggedleft \hfill (Chapter~\ref{cha:soa})}
    \begin{itemize}
        \item Comprehensive review of hardware accelerators for sparse matrix computing.
        \item Identification of common architectural features and design tendencies.
        \item Comparative analysis to explain design choices under varying constraints.
    \end{itemize}

    \item \textbf{Design of a cache architecture for ``random'' \glspl{reduction}} {\raggedleft \hfill (Chapter~\ref{cha:spdc})}
    \begin{itemize}
        \item Proposal of a cache with native \gls{reduction} support, capable of handling both dense and sparse data via multiple cache regions.
        \item Development of a Python-based simulation model to analyze memory transactions and benchmark against a generic cache.
    \end{itemize}

    \item \textbf{Probabilistic modeling and exploration of cache behavior} {\raggedleft \hfill (Chapter~\ref{cha:theory})}
    \begin{itemize}
        \item Implementation of a lightweight C-based probabilistic model for rapid simulation.
        \item Exploration of cache dimensioning and behavior across diverse workloads.
    \end{itemize}

    \item \textbf{Micro-architecture and evaluation through \acrshort{rtl} simulation} {\raggedleft \hfill (Chapter~\ref{cha:arch})}
    \begin{itemize}
        \item Introduction of a memory partitioning scheme enabling multiple configurable regions.
        \item Assessment of the proposed cache architecture using \acrshort{rtl}-level simulation.
        \item Identification of two groups of sparse matrices leading to different performance profiles.
    \end{itemize}

    \item \textbf{Future directions for performance and scalability} {\raggedleft \hfill (Chapter~\ref{cha:optim})}
    \begin{itemize}
        \item Extension of the cache architecture to multicore systems for improved scalability.
        \item Optimization techniques to enhance usability and practical deployment.
    \end{itemize}
\end{enumerate}

\cleardoublepage %% Comment this line if this section occupies the last pages